\documentclass[11pt]{article}
\usepackage{graphicx} % Required for inserting images
\usepackage[margin=2.5cm]{geometry}
\usepackage{listings}
\usepackage{array}
\usepackage{booktabs}
\usepackage{xcolor}
\usepackage{hyperref}

\begin{document}

\begin{center}
    \Large{\textsc{Carta del Uno}}
\end{center}

\vspace{2px}

\noindent\textbf{Para}: Moreno López, $\forall$malia.

\vspace{2px}

\noindent\textbf{De}: mío.

\vspace{2px}

\noindent\textbf{Sujecto}: elíptico.

\vspace{2px}

\noindent\textbf{Fecha}: 04--04 (no pongo el año y así tengo la del año que viene as well).

\vspace{6px}

\hrule

\vspace{6px}

\noindent Para empezar por donde cabría esperar, felicitaciones por cumplir años. Ya son 20, un hito reseñable. Quizás esté pareciendo un
inicio poco serio, pero aseguro que no va a seguir así. \\

\noindent En febrero, cuando empezamos a quedar de nuevo\footnote{Tras unos 2 años de barbecho.}, fue muy grato recordar lo extremadamente agradable que es \textit{tu compañía}. Desde mi punto de vista,
pasar tiempo contigo es siempre una experiencia que desearía no terminara. Prácticamente cada quedada desde que nos conocemos, que no es poco, tiene reservada un hueco especial en mi \textit{memoria episódica}. 

No sé cuál será el motivo, a veces parece que posees habilidades sobrehumanas, pero trasmites una \textit{paz} y una \textit{alegría} 
desmedidas\footnote{Pondría la mano en el fuego porque mucha gente piensa igual.}, capaces de mejorar el día de cualquiera. En cierta manera, me hace poner en duda el extendido dicho

\begin{center}
    \textit{``Recibes lo que das''},
\end{center}

\noindent pues poca gente puede dar tanto como lo haces tú. \\

\noindent Por otro lado, cambiando de temática, me gustaría hablar brevemente sobre la admiración que siento hacia tu persona y forma de ser. No miento\footnote{Ya sabes tú que a mí no me van vuestras costumbres,} si digo
que eres, sin lugar a dudas, una de las personas más \textit{extraordinarias} que conozco, pues cada vez que pienso en alguien de dichas características, tú eres casi siempre la primera persona que me viene a la mente. % No solo por tu capacidad de \textit{empatizar} con los demás, sino también por tu \textit{inteligencia} y \textit{sabiduría}. \\

Ciertamente, no es fácil encontrar a alguien que sepa \textit{escuchar}\footnote{Sobre todo cuando tus oídos tienen que soportar aberraciones como Xavibo.} y \textit{comprender} a los demás tan pacientemente como lo haces tú, dando por descontado tu más que sobresaliente inteligencia. \\ % No me malinterpretes, no quiero decir que seas la única persona capaz de hacerlo, pero sí que eres la mejor. \\

\noindent Para terminar, incido en que eres bastante especial y, por otra parte, que espero que la redacción sea de tu agrado, es lo más serio que he hecho desde los comentarios críticos, aunque quizás no debería preocuparme sabiendo a la lírica que estás acostumbrada. \\

\noindent Atentamente,

\texttt{\@@samuufdzz\_}

\end{document}
